\chapter{Experiments and Evaluation}
\label{ch:experiments}

% Generally: Each experiment should have its own question, setup, results, and interpretation.
% what follows is just an idea. Don't take too rigidly
% need to prove how purely waveform prediction - cost and cost fails
% --------


% Chapter introduction:
% - Explain how the experiments answer the design questions formulated in Chapter 4.
% - Establish the sequence: select geometry and training settings, compare Predictors
%   and numerical methods, then evaluate closed-loop stimulation.
% - End each study with the evidence supporting the choices used in subsequent studies.
% Editorial boundary: refer to Chapter 4 for method definitions. Give experimental
% settings, search spaces/budgets, results, and interpretation here.
% For each comparison, state the question, what varies, what is held fixed, and the
% criterion for judging the outcome. Do not assume that the expected improvement occurs.

\section{Experimental Setup and Evaluation Criteria}

\subsection{Simulation Scenarios and Data Splits}
% Specify Plant configurations, simulation durations, stimulation scenarios, dataset
% sizes, and independent repetitions. Refer to Chapter 4 for construction of the Plant.
% Describe healthy and unstimulated seizure activity as reference conditions.
% Include evidence for Plant calibration here if that calibration is an empirical study.
% State which trajectories/seeds are used for training, selection, and final evaluation.
% Use matched conditions across methods where possible and report variability across runs.
% QUESTION: Which seizure regimes and stimulation conditions must the conclusions cover?

\subsection{Prediction Accuracy and Seizure Metrics}
% Define evaluation measures and aggregation over channels, time, and repetitions.
% Include waveform and spectral prediction errors, Seizure State, and Seizure Burden
% as relevant. Distinguish sensor-space scores from regional ground truth.
% Assess prediction error over physical forecast duration, not just the number of steps.
% Keep training Loss, controller Cost, and Closed-Loop Evaluation Score distinct.
% QUESTION: What operational criterion defines a useful prediction duration for control?

\subsection{Stimulation Use and Computational Performance}
% Define measures for current magnitude, stimulation duration, distinct pulse count,
% and constraint satisfaction according to the formulations in Chapter 4.
% State the computational setup and how solve time, failures, and solution quality are measured.
% QUESTION: Is a controller update deadline part of the study, or is runtime a comparison
% measure only? Do not claim real-time feasibility without testing a stated deadline.

\section{Selection of Spectral Observable Geometry}
% QUESTION: Which STFT geometry makes the spectral Observable useful for prediction
% and control, under the selection criterion established in Chapter 4?

\subsection{Geometry Search and Comparison Protocol}
% Give candidate Segment lengths, Frame hops, and other geometry parameters varied.
% State the Predictor class, fitting procedure, search budget, and selection data.
% Account for changes in target dimension, Frame rate, and physical prediction duration.
% QUESTION: What common evaluation reference makes scores comparable when changing
% the geometry changes the target itself? Avoid ranking solely by native training Loss.

\subsection{Prediction Accuracy and Temporal Responsiveness}
% Compare prediction quality and preservation of seizure onset/propagation dynamics.
% Report tradeoffs associated with temporal and spectral resolution and smoothing.
% If selection uses closed-loop performance, include that evidence on selection scenarios
% here and reserve independent scenarios for the final closed-loop comparison.
% QUESTION: Does a geometry that is easier to predict still resolve changes needed for control?

\subsection{Selected Geometry}
% Justify the configuration carried into later experiments using the declared criterion.
% Report sensitivity to nearby choices and variability rather than only the best trial.
% State what the evidence establishes about this choice and where it remains uncertain.

\section{Waveform Predictor Training-Loss Comparison}
% QUESTION: How do waveform Loss, spectral Loss, and curriculum training affect the
% accuracy and useful duration of waveform predictions?

\subsection{Loss Geometry and Training Configurations}
% Compare the Loss combinations and curricula defined in Chapter 4, including the
% intended comparison of STFT training with curriculum plus STFT training.
% Report searches over spectral Loss geometry, weights, and Spans where investigated.
% This geometry scores waveform predictions; it need not equal the spectral Observable geometry.
% Control model capacity, data, training effort, and repetitions, or explain differences.
% QUESTION: Which comparisons isolate the effect of the Loss, geometry, and curriculum
% rather than changing all three at once?

\subsection{Prediction Results and Training Strategy Selection}
% Evaluate waveform and spectral errors on common criteria across training objectives.
% Examine multi-step error and retention of seizure dynamics and stimulation responses.
% Justify the Loss geometry and training strategy used in the Predictor comparison.
% Assess their consequences for control in the closed-loop study below; if control is
% the selection criterion, report selection evidence here separately from final evaluation.

\section{Predictor Comparison}
% QUESTION: What is gained or lost by predicting waveforms versus spectral Observables,
% and by using linear/DMD versus nonlinear MLP dynamics?

\subsection{Comparison Protocol and Model Selection}
% List the representation/model combinations studied and the configurations selected above.
% Report remaining hyperparameter searches and comparable selection budgets.
% Separate the representation comparison from the model-class comparison where possible.
% QUESTION: Are the selected geometry and Loss settings transferable across model classes,
% or does each class require its own selection? State the fairness criterion used.

\subsection{Prediction Quality over the Control Horizon}
% Compare errors at matched physical forecast durations and across relevant seizure regimes.
% Use shared spectral evaluation criteria where waveform errors cannot be compared directly.
% Assess response to candidate stimulation, as well as unstimulated seizure evolution.
% Identify Predictors carried into closed-loop evaluation without equating prediction
% accuracy with demonstrated control performance.

\section{Numerical Solution Comparison}
% QUESTION: Which formulation and solver solve the control problems reliably at the
% Predictor dimensions and Control Horizons used in the study?
% Keep this section proportional to its scientific role; shorten it if solver selection
% is only a supporting choice. Definitions belong in Chapter 4.

\subsection{Single Shooting, Multiple Shooting, and PCMS}
% Compare the selected transcriptions, including the NARX and PCMS formulations.
% Hold the underlying Predictor, Cost, constraints, and initial control problem fixed.
% State initialization, termination tolerances, and treatment of whole-horizon Costs.

\subsection{Solution Quality, Reliability, and Runtime}
% Compare attained Cost, constraint violations, failure frequency, and solve time.
% Examine the effect of Predictor class and Control Horizon where relevant.
% Justify the numerical methods used in the subsequent closed-loop experiments.
% QUESTION: Does a faster solve preserve the objective and feasibility of the original problem?

\section{Closed-Loop Control Evaluation}
% Test whether the choices above lead to seizure suppression under admissible stimulation.
% Use final evaluation scenarios that were not used to select configurations.

\subsection{Evaluation Protocol and Reference Conditions}
% Specify controller configurations, Control Horizons, update intervals, and evaluation durations.
% Include the unstimulated reference and justify any additional control baselines.
% State how controller Cost weights and budgets were selected before final evaluation.
% QUESTION: Which comparisons isolate the Predictor, suppression Cost, and stimulation
% formulation, and which compare the best selected combination for each method?

\subsection{Prediction Representation and Suppression Objective}
% QUESTION: Do the open-loop prediction advantages translate into lower Seizure Burden?
% Compare waveform and spectral Predictors and the suppression Costs defined in Chapter 4.
% Relate sensor-space improvement to regional seizure outcomes and stimulation use.
% Include focused training-Loss comparisons if needed to establish their effect on control.

\subsection{Continuous and Pulse-Based Stimulation}
% QUESTION: How do L2/L1 continuous control and pulse-based formulations trade seizure
% suppression against current magnitude, stimulation time, pulse count, and computation?
% Compare the mixed-integer method, relaxation with schedule recovery, and continuous
% pulse-timing method included in the thesis, under explicitly comparable restrictions.
% Evaluate recovered schedules as delivered to the Plant, not just their relaxed objectives.
% QUESTION: What changes after schedule recovery in feasibility and achieved suppression?
% QUESTION: Does optimizing amplitude or duty cycle add benefit over the simpler pulse formulation?

\subsection{Tradeoffs and Limits of the Findings}
% Bring together suppression, stimulation use, and computational reliability across runs.
% Identify where prediction-based selection agrees with closed-loop performance and where it fails.
% Distinguish conclusions supported across scenarios from those specific to the simulated setup.
% Close by answering the chapter's questions and preparing the thesis conclusion.
